\section{Conclusion}

\projectname{} was developed to address a clear need in microbiome informatics: the need to centralize literature-derived evidence in a form that preserves study context and remains usable for later exploration. Rather than treating papers as isolated notes or reducing results to flat organism lists, the project models studies, groups, comparisons, taxa, and findings explicitly. This makes the platform more suitable for structured querying, reproducible curation, and future synthesis across publications.

The current implementation demonstrates that this approach is practical. A Django and PostgreSQL architecture provides a stable backbone for data management, the web interface supports browsing and exploration, the import workflow adds validation and provenance, and Taxon Bridge strengthens taxonomy-aware normalization. Although the system is still at prototype stage, it already provides a credible foundation for centralized microbiome evidence integration and future bioinformatics work built on curated literature data.
